\section{Proactive: la diferencia esencial entre los productos de IA y los productos de internet}

Tristan considera que el mayor cambio de interacción en la era de la IA no es LUI ni GUI, sino lo proactive: por fin hay algo que puede actuar por iniciativa propia en tu nombre. OpenClaw permitió experimentar lo proactive en la capa de herramientas; Elys, en cambio, lleva lo proactive a la capa social y hace que el doble digital salga por iniciativa propia a relacionarse con el mundo.

\begin{figure}[H]
\centering
\includegraphics[width=0.92\textwidth]{proactive-product-shift.png}
\caption{Paradigmas de producto: de Reactive a Assistive y de ahí a Proactive.}
\end{figure}

\begin{knowledgebox}{Lectura de la figura: Proactive cambia el modelo mental del producto}
Los productos Reactive requieren que el usuario tome la iniciativa; los productos Assistive ayudan mientras el usuario opera; los productos Proactive pueden explorar oportunidades en nombre del usuario cuando este no hace nada por iniciativa propia. El potencial de Elys proviene del tercer nivel: el doble digital socializa por iniciativa propia y le trae conexiones al usuario.
\end{knowledgebox}

\subsection{La competencia con las empresas de modelos}

Tristan considera que la IA social no necesariamente requiere entrenar un modelo base propio. La razón es que el núcleo de una red social no es solo la inteligencia, sino el context, el paradigma de producto y la conexión entre personas reales. Las empresas de modelos se enfocan en la inteligencia en sí y no necesariamente tienen lo social como primera prioridad; si una empresa de aplicaciones logra controlar el volante de inercia del context, también puede tener una oportunidad independiente.

\begin{importantbox}{¿Dónde está la barrera de entrada de una empresa de aplicaciones?}
No se trata de «yo también puedo llamar al modelo», sino de si tengo un context propio, relaciones reales, un modelo mental de producto, retroalimentación de datos y la confianza de los usuarios. Sin esto, las empresas de modelos absorberán la aplicación; con esto, la aplicación puede generar sus propios efectos de red.
\end{importantbox}

\subsection{Los seres humanos están demasiado solos}

Al final, la entrevista vuelve al núcleo espiritual: los seres humanos están demasiado solos. Eve intenta resolver la soledad mediante la compañía de la IA; Elys intenta resolverla mediante la IA como promotora de conexiones entre personas reales. Lo primero es una relación entre la IA y la persona; lo segundo es la IA ayudando a que las personas establezcan relaciones entre sí. Esta diferencia determina que Elys no sea un producto de compañía, sino un producto de red social.

\begin{warningbox}{El valor emocional no es una necesidad menor}
Socializar, tener compañía, sentirse comprendido y conectar son necesidades humanas frecuentes y profundas. Si un producto de IA logra atender mejor estas necesidades de forma segura, auténtica y controlable, puede estar más cerca del usuario común que un producto puramente utilitario. Pero el riesgo también es mayor, porque toca la privacidad, la identidad y las relaciones.\end{warningbox}

\subsection{Resumen de la sección}

Lo proactive es lo que distingue a los productos de IA de los productos de internet. El núcleo de Elys no es que la IA converse para acompañar, sino que el doble digital de IA reduzca por iniciativa propia la fricción social en nombre del usuario y, en última instancia, le traiga conexiones con personas reales.

